\documentclass[11pt]{article}
\usepackage[margin=1in]{geometry}
\usepackage{amsmath,amssymb,amsthm}
\usepackage[utf8]{inputenc}
\usepackage{listings}
\usepackage{hyperref}
\usepackage{xcolor}
\usepackage{textcomp}

\definecolor{leanblue}{rgb}{0.2,0.4,0.7}

\lstdefinelanguage{Lean}{
  keywords={theorem,lemma,def,axiom,structure,class,instance,where,by,sorry,simp,exact,intro,apply,have,let,fun,match,if,then,else,import,open,namespace,end,section,variable,inductive,noncomputable,abbrev,deriving},
  keywordstyle=\color{leanblue}\bfseries,
  comment=[l]{--},
  morecomment=[s]{/-}{-/},
  string=[b]",
  basicstyle=\ttfamily\small,
  breaklines=true,
  extendedchars=true,
  inputencoding=utf8,
  literate=
    {≥}{{$\geq$}}1
    {≤}{{$\leq$}}1
    {→}{{$\to$}}1
    {←}{{$\leftarrow$}}1
    {∧}{{$\wedge$}}1
    {∨}{{$\vee$}}1
    {¬}{{$\neg$}}1
    {∀}{{$\forall$}}1
    {∃}{{$\exists$}}1
    {⊆}{{$\subseteq$}}1
    {⊂}{{$\subset$}}1
    {∈}{{$\in$}}1
    {∉}{{$\notin$}}1
    {×}{{$\times$}}1
    {⊗}{{$\otimes$}}1
    {▸}{{$\blacktriangleright$}}1
    {⟨}{{$\langle$}}1
    {⟩}{{$\rangle$}}1
    {λ}{{$\lambda$}}1
    {α}{{$\alpha$}}1
    {β}{{$\beta$}}1
    {σ}{{$\sigma$}}1
    {θ}{{$\theta$}}1
    {ε}{{$\varepsilon$}}1
    {μ}{{$\mu$}}1
    {π}{{$\pi$}}1
    {Σ}{{$\Sigma$}}1
    {Π}{{$\Pi$}}1
    {▶}{{$\triangleright$}}1
    {⊔}{{$\sqcup$}}1
    {⊓}{{$\sqcap$}}1
    {⊥}{{$\bot$}}1
    {⊤}{{$\top$}}1
    {∅}{{$\emptyset$}}1
    {∪}{{$\cup$}}1
    {∩}{{$\cap$}}1
    {≠}{{$\neq$}}1
    {ℕ}{{$\mathbb{N}$}}1
    {₀}{{$_0$}}1
    {₁}{{$_1$}}1
    {₂}{{$_2$}}1
    {|>}{{$|{\!>}$}}2
    {·}{{$\cdot$}}1
    {—}{{---}}1
    {∘}{{$\circ$}}1
    {√}{{$\sqrt{}$}}1
    {∝}{{$\propto$}}1
    {≈}{{$\approx$}}1
    {═}{{=}}1
    {⟹}{{$\Longrightarrow$}}1,
}

\newtheorem{theorem}{Theorem}
\newtheorem{lemma}[theorem]{Lemma}
\newtheorem{definition}[theorem]{Definition}
\newtheorem{axiom_stmt}[theorem]{Axiom}

\title{Formal Verification: Vouch-Based Trust Model}
\author{Zach Kelling, Lux Industries}
\date{December 2025}

\begin{document}
\maketitle

\begin{abstract}
We formalize the vouch-based trust model where trust flows through attestation chains without any global root or certificate authority. Each vouch bounds the scope of delegated authority, and chain authority is monotonically decreasing.
\end{abstract}

\section{Introduction}
The vouch-based trust model replaces traditional PKI with attestation chains. No CA, no global root. Trust flows through signed vouches, each bounding the scope of delegated authority. This model is used by Hanzo IAM, Lux validator trust, and Zoo role endorsements.

\section{Definitions}
\begin{definition}[RecognitionLevel]
Trust levels: unrecognized (0), vouched (1), direct (2).
\end{definition}

\begin{definition}[Vouch]
A vouch: voucher, vouchee, scope bound, expiry timestamp.
\end{definition}


\section{Theorems and Axioms}
\begin{theorem}[\texttt{unrecognized\_min}]
Unrecognized is the minimum recognition level.
\end{theorem}

\begin{theorem}[\texttt{chain\_authority\_monotone}]
Adding a vouch to a chain can only reduce authority: $\text{chainAuthority}(v :: \text{chain}, \text{top}) \leq \text{chainAuthority}(\text{chain}, \text{top})$.
\end{theorem}
\begin{proof}
By induction on the chain, using $\min$ monotonicity.
\end{proof}

\begin{theorem}[\texttt{empty\_chain\_full}]
Empty chain gives full authority: $\text{chainAuthority}([], \text{top}) = \text{top}$.
\end{theorem}

\begin{theorem}[\texttt{chain\_bounded}]
Chain authority is bounded by the top: $\text{chainAuthority}(\text{chain}, \text{top}) \leq \text{top}$.
\end{theorem}

\begin{theorem}[\texttt{chain\_bounded\_by\_scope}]
Chain authority is bounded by every vouch's scope in the chain.
\end{theorem}
\begin{proof}
By induction: each vouch's scope constrains the $\min$ fold.
\end{proof}


\section{Lean Source}
\lstinputlisting[language=Lean]{../../formal/lean/Trust/Vouch.lean}

\section{References}
\noindent Source code: \url{https://github.com/luxfi/proofs/blob/main/lean/Trust/Vouch.lean}\\
\noindent White papers: \url{https://papers.lux.network}

\noindent Related white papers:
\begin{itemize}
  \item \texttt{lux-id-iam.tex}
  \item \texttt{lux-id-did-specification.tex}
\end{itemize}

\end{document}
